% 第1章: 全部main.cppに書くとどうなる？
\section{全部main.cppに書くとどうなる？}

\begin{notebox}[C++クラスの復習が必要な方へ]
この章からC++のクラスを使います。構造体・クラス・コンストラクタに自信がある方はそのまま読み進めてください。不安な方は第1章「前提知識 — C++クラスの基礎」（\pageref{sec:prerequisite}ページ）で復習できます。
\end{notebox}

\begin{cautionbox}[サンプルコードについて]
本リポジトリに実装されているモジュール群のコードは、あくまで\textbf{設計パターンを具体的に示すためのサンプル実装}である。
実機での動作検証は行っていないため、実プロジェクトへの適用時は十分なテストを実施すること。
\end{cautionbox}

%% ── Q1: ベタ書きの問題 ──────────────────────
\begin{qabox}{全部main.cppに書くの、何がまずいの？}

まずは「よくある書き方」を見てみましょう。LED2個と温度センサー1個を制御するプログラムです。

\begin{badexample}{全部main.cppにベタ書き}
\begin{lstlisting}
#include <Arduino.h>

// ピン番号（マジックナンバー）
#define LED1_PIN 2
#define LED2_PIN 4
#define TEMP_PIN 34

// グローバル変数がバラバラに並ぶ
bool led1State = false;
bool led2State = false;
float temperature = 0.0;
unsigned long lastBlink = 0;
unsigned long lastRead = 0;

void setup() {
    Serial.begin(115200);
    pinMode(LED1_PIN, OUTPUT);
    pinMode(LED2_PIN, OUTPUT);
    // 温度センサーはアナログなのでpinMode不要
}

void loop() {
    unsigned long now = millis();

    // LED1: 500msごとに点滅
    if (now - lastBlink >= 500) {
        lastBlink = now;
        led1State = !led1State;
        digitalWrite(LED1_PIN, led1State ? HIGH : LOW);
    }

    // LED2: 常時点灯
    digitalWrite(LED2_PIN, HIGH);

    // 温度センサー: 1秒ごとに読み取り
    if (now - lastRead >= 1000) {
        lastRead = now;
        int raw = analogRead(TEMP_PIN);
        temperature = raw * 3.3 / 4095.0 * 100.0;
        Serial.print("Temp: ");
        Serial.println(temperature);
    }

    // ここにモーター制御を追加したら？
    // さらにディスプレイ表示を追加したら？
    // ...どんどん膨れ上がる
}
\end{lstlisting}
\end{badexample}

動くことは動きます。でも、この書き方には3つの大きな問題があります。

\subsection{問題1: 肥大化}

今は約40行ですが、モーター制御、ディスプレイ表示、通信機能……と追加していくと、あっという間に数百行になります。「LED1の点滅ロジックはどこだっけ？」と探すだけで一苦労です。

\subsection{問題2: マジックナンバー}

\texttt{LED1\_PIN}、\texttt{500}（点滅間隔）、\texttt{3.3 / 4095.0 * 100.0}（変換式）といった数値がコード中に散らばっています。後から見て「この数字は何？」となりがちです。

\subsection{問題3: 再利用できない}

別のプロジェクトでもLED点滅が必要になったとき、このコードからLED関連の部分だけを抜き出すのは大変です。変数名が被ったり、依存関係が絡み合っていたりします。

\end{qabox}

%% ── Q2: 教科書の構成 ──────────────────────
\begin{qabox}{この教科書で何を学ぶの？}

この教科書では、上のような「ベタ書きコード」を段階的に改善していきます。

\begin{enumerate}
  \item \textbf{第3章: モジュールという考え方} — 1つのハードウェアを1つのクラスに分離する
  \item \textbf{第4章: インターフェースで統一する} — 異なるモジュールを同じ方法で扱う
  \item \textbf{第5章: データの流れを整理する} — モジュール間のデータ受け渡しを安全にする
  \item \textbf{第6章: 3フェーズモデル} — loop()の処理を「入力→ロジック→出力」に整理する
  \item \textbf{第7章: プロジェクト構成を理解する} — フォルダ配置とPlatformIOの使い方
  \item \textbf{第8章: 実践} — 温度センサーモジュールを一から作る
  \item \textbf{第9章: 設計のコツ集} — 実務で役立つTips
\end{enumerate}

最終的に、先ほどのベタ書きコードがどう変わるか？ それは第6章で明らかになります。

\end{qabox}
